\documentclass[11pt]{article}
\usepackage[margin=1in]{geometry}
\usepackage{amsmath, amssymb, amsthm}
\usepackage{enumitem}
\usepackage{xcolor}
\usepackage{tcolorbox}
\usepackage{graphicx}
\tcbuselibrary{skins, breakable}

% Page setup
\pagestyle{plain}

% Color definitions
\definecolor{problemconfig}{RGB}{230, 242, 255}
\definecolor{strategic}{RGB}{255, 230, 242}
\definecolor{tactical}{RGB}{242, 255, 230}
\definecolor{techniques}{RGB}{255, 242, 230}
\definecolor{verification}{RGB}{255, 230, 230}
\definecolor{solution}{RGB}{230, 255, 255}
\definecolor{competition}{RGB}{242, 242, 255}
\definecolor{gold}{RGB}{218, 165, 32}

% Box styles
\newtcolorbox{problemconfigbox}{
    colback=problemconfig,
    colframe=blue,
    title=Problem Configuration Analysis,
    fonttitle=\bfseries,
    arc=3pt,
    boxrule=1pt,
    breakable
}

\newtcolorbox{strategicbox}{
    colback=strategic,
    colframe=purple,
    title=Strategic Framework Application,
    fonttitle=\bfseries,
    arc=3pt,
    boxrule=1pt,
    breakable
}

\newtcolorbox{tacticalbox}{
    colback=tactical,
    colframe=green,
    title=Tactical Methodology Selection,
    fonttitle=\bfseries,
    arc=3pt,
    boxrule=1pt,
    breakable
}

\newtcolorbox{techniquesbox}{
    colback=techniques,
    colframe=orange,
    title=Conversion Techniques Application,
    fonttitle=\bfseries,
    arc=3pt,
    boxrule=1pt,
    breakable
}

\newtcolorbox{verificationbox}{
    colback=verification,
    colframe=red,
    title=Verification Strategy Design,
    fonttitle=\bfseries,
    arc=3pt,
    boxrule=1pt,
    breakable
}

\newtcolorbox{solutionbox}{
    colback=solution,
    colframe=cyan,
    title=Solution Roadmap,
    fonttitle=\bfseries,
    arc=3pt,
    boxrule=1pt,
    breakable
}

\newtcolorbox{competitionbox}{
    colback=competition,
    colframe=purple,
    title=Competition Strategy Integration,
    fonttitle=\bfseries,
    arc=3pt,
    boxrule=1pt,
    breakable
}

\newtcolorbox{keyinsightbox}{
    colback=yellow!20,
    colframe=gold,
    title=Key Insight,
    fonttitle=\bfseries,
    arc=3pt,
    boxrule=2pt,
    leftrule=5pt,
    breakable
}

\newtcolorbox{warningbox}{
    colback=red!10,
    colframe=red!50,
    title=Warning: Common Pitfall,
    fonttitle=\bfseries,
    arc=3pt,
    boxrule=2pt,
    leftrule=5pt,
    breakable
}

\newtcolorbox{tipbox}{
    colback=green!10,
    colframe=green!50,
    title=Pro Tip,
    fonttitle=\bfseries,
    arc=3pt,
    boxrule=2pt,
    leftrule=5pt,
    breakable
}

\title{\textbf{2017 AMC 12B Problem 15: Strategic Analysis}}
\author{Using AMC12/COMC Problem-Solving Framework}
\date{\today}

\begin{document}
\maketitle

\section*{Problem Statement}
Let $ABC$ be an equilateral triangle. Extend side $\overline{AB}$ beyond $B$ to a point $B'$ so that $BB'=3AB$. Similarly, extend side $\overline{BC}$ beyond $C$ to a point $C'$ so that $CC'=3BC$, and extend side $\overline{CA}$ beyond $A$ to a point $A'$ so that $AA'=3CA$. What is the ratio of the area of $\triangle A'B'C'$ to the area of $\triangle ABC$?

\textbf{Answer Choices:} 
$\textbf{(A)}\ 9:1 \qquad\textbf{(B)}\ 16:1 \qquad\textbf{(C)}\ 25:1 \qquad\textbf{(D)}\ 36:1 \qquad\textbf{(E)}\ 37:1$

\begin{problemconfigbox}
\subsection*{A. Problem Classification}
\begin{itemize}[leftmargin=*]
    \item \textbf{Problem Type}: To Find - determine a specific numerical ratio
    \item \textbf{Difficulty Band}: Mid-band (Problem 15 of 25) - requires geometric insight and calculation
    \item \textbf{Competition Context}: AMC12 (75 min, 25 problems) - approximately 3 minutes allocated
    \item \textbf{Mathematical Domain}: Geometry (triangle areas, symmetry, similar figures)
\end{itemize}

\subsection*{B. Problem Structure Identification}
\begin{itemize}[leftmargin=*]
    \item \textbf{Given Information}: 
    \begin{itemize}
        \item $ABC$ is equilateral (all sides equal, all angles $60^{\circ}$)
        \item Extensions: $BB' = 3AB$, $CC' = 3BC$, $AA' = 3CA$
        \item Symmetric construction pattern
    \end{itemize}
    \item \textbf{Constraints}: 
    \begin{itemize}
        \item Extensions are along existing sides (collinear)
        \item All extensions have the same ratio (3:1)
        \item Must maintain geometric relationships
    \end{itemize}
    \item \textbf{Objective}: Find ratio $\frac{[A'B'C']}{[ABC]}$
    \item \textbf{Hidden Assumptions}: 
    \begin{itemize}
        \item $\triangle A'B'C'$ is also equilateral by symmetry
        \item Can decompose $\triangle A'B'C'$ into smaller regions
        \item Area ratios relate to side length ratios squared
    \end{itemize}
    \item \textbf{Answer Choice Leverage}: 
    \begin{itemize}
        \item Options (A)-(D) are perfect squares suggesting integer side length ratio
        \item Option (E) $= 37 = 36 + 1$ hints at decomposition strategy
        \item Can eliminate unreasonable bounds quickly
    \end{itemize}
\end{itemize}

\subsection*{C. Strategic Orientation}
\begin{itemize}[leftmargin=*]
    \item \textbf{Hypothesis}: By symmetry, $\triangle A'B'C'$ is equilateral
    \item \textbf{Conclusion}: Need to find either side length ratio or direct area ratio
    \item \textbf{Notation}: Let $AB = BC = CA = s$ (side length of original triangle)
    \item \textbf{Intuitive Approach}: 
    \begin{itemize}
        \item Decompose large triangle into original plus three corner triangles
        \item Use Law of Cosines to find side of $\triangle A'B'C'$
        \item Apply area formula or coordinate geometry
    \end{itemize}
    \item \textbf{Cross-Domain Potential}: 
    \begin{itemize}
        \item Pure geometry (Law of Cosines, area formulas)
        \item Coordinate geometry (place triangle on coordinate system)
        \item Vector methods (barycentric coordinates)
    \end{itemize}
\end{itemize}
\end{problemconfigbox}

\begin{strategicbox}
\subsection*{A. Psychological Strategies}
\begin{itemize}[leftmargin=*]
    \item \textbf{Mental Toughness}: Multiple solution paths exist; if one fails, try another approach
    \item \textbf{Creativity Cultivation}: Recognize symmetry immediately; visualize decomposition
    \item \textbf{Iterative Refinement}: Start with simplest approach (side length), escalate if needed
\end{itemize}

\subsection*{B. Investigation Strategies}
\begin{itemize}[leftmargin=*]
    \item \textbf{Startup Orientation}: Draw accurate diagram, label all points, identify symmetry
    \item \textbf{Get Your Hands Dirty}: Set $AB = 1$ or $AB = 2$ for concrete calculations
    \item \textbf{Penultimate Step}: Once side of $\triangle A'B'C'$ is found, square it for area ratio
    \item \textbf{Wishful Thinking}: Wish for simple integer or radical answer (guides verification)
    \item \textbf{Make It Easier}: Consider just one corner triangle first, then multiply by 3
    \item \textbf{Conjecture via Small Cases}: Check answer choices against extreme cases
    \item \textbf{Visualization}: Draw clear diagram showing extensions and resulting triangles
\end{itemize}

\subsection*{C. Late-Band Strategic Playbook}
\begin{itemize}[leftmargin=*]
    \item \textbf{Answer-Choice Leverage}: Option (E) being $37 = 36 + 1$ suggests $36 = 3 \times 12$ decomposition
    \item \textbf{Make It Easier}: Focus on finding one side of $\triangle A'B'C'$ instead of full area
    \item \textbf{Penultimate Step}: Area ratio = (side ratio)$^2$ for similar triangles
    \item \textbf{Elimination}: Quick bounds check rules out (A) and (B) as too small
\end{itemize}

\subsection*{D. Argument Construction Strategy}
\begin{itemize}[leftmargin=*]
    \item \textbf{Direct Proof}: Calculate side of $\triangle A'B'C'$ using Law of Cosines, then square
    \item \textbf{Decomposition Argument}: Split $\triangle A'B'C'$ into four triangles, sum areas
    \item \textbf{WLOG}: Without loss of generality, let $AB = 1$ or $AB = 2$ (scales cancel)
\end{itemize}

\subsection*{E. Cross-Domain Translation}
\begin{itemize}[leftmargin=*]
    \item \textbf{Draw a Picture}: Essential - draw triangle with extensions to visualize structure
    \item \textbf{Recast the Problem}: Transform to coordinate geometry if pure geometry stalls
    \item \textbf{Change Point of View}: View as three congruent triangles surrounding original
\end{itemize}
\end{strategicbox}

\begin{tacticalbox}
\subsection*{A. Core Mathematical Tactics}
\begin{itemize}[leftmargin=*]
    \item \textbf{Symmetry Exploitation}: 
    \begin{itemize}
        \item All three extensions are identical by rotational symmetry
        \item $\triangle A'B'C'$ must be equilateral
        \item Three corner triangles are congruent
    \end{itemize}
    \item \textbf{Wishful Thinking}: Assume answer has clean form, guides calculation strategy
    \item \textbf{Normalization}: WLOG set $AB = 1$ or $AB = 2$ to simplify arithmetic
    \item \textbf{Decomposition}: Break complex region into manageable pieces
    \item \textbf{Invariant Principle}: Area ratios independent of actual side length (scaling)
\end{itemize}

\subsection*{B. Domain-Specific Tactics}
\begin{itemize}[leftmargin=*]
    \item \textbf{Geometry}: 
    \begin{itemize}
        \item Law of Cosines for finding side lengths
        \item Area formula: $A = \frac{1}{2}ab\sin\theta$
        \item Equilateral triangle area: $A = \frac{s^2\sqrt{3}}{4}$
        \item Similar triangles: area ratio = (side ratio)$^2$
    \end{itemize}
    \item \textbf{Coordinate Geometry}: Place triangle strategically (e.g., $B$ at origin)
    \item \textbf{Vector Methods}: Barycentric coordinates for advanced approach
\end{itemize}
\end{tacticalbox}

\begin{techniquesbox}
\subsection*{A. Recognition Index}
\begin{itemize}[leftmargin=*]
    \item \textbf{Pattern Recognition}: 
    \begin{itemize}
        \item Symmetric construction $\Rightarrow$ symmetric result
        \item Equilateral triangle $\Rightarrow$ all angles $60^{\circ}$ or $120^{\circ}$
        \item Extension ratio 3:1 $\Rightarrow$ sides become 4:1
    \end{itemize}
    \item \textbf{Technique Triggers}:
    \begin{itemize}
        \item ``Extend beyond'' $\Rightarrow$ Law of Cosines likely useful
        \item ``Ratio of areas'' $\Rightarrow$ Consider both direct and decomposition methods
        \item Equilateral triangle $\Rightarrow$ Exploit $60^{\circ}$ angles
    \end{itemize}
\end{itemize}

\subsection*{B. Technique Selection}
\begin{itemize}[leftmargin=*]
    \item \textbf{Primary Technique}: Law of Cosines
    \begin{itemize}
        \item To find $B'C'$: Given $BB' = 3s$, $BC' = BC + CC' = 4s$
        \item Angle $\angle B'BC' = 180^{\circ} - 60^{\circ} = 120^{\circ}$
        \item Apply: $(B'C')^2 = (BB')^2 + (BC')^2 - 2(BB')(BC')\cos(120^{\circ})$
    \end{itemize}
    \item \textbf{Alternative Technique}: Area Decomposition
    \begin{itemize}
        \item $[A'B'C'] = [ABC] + [AA'B'] + [BB'C'] + [CC'A']$
        \item Each corner triangle has area using $\frac{1}{2}ab\sin\theta$
        \item $[AA'B'] = \frac{1}{2} \cdot AA' \cdot AB' \cdot \sin(120^{\circ})$
    \end{itemize}
\end{itemize}

\subsection*{C. Technique Integration}
\begin{itemize}[leftmargin=*]
    \item \textbf{Workflow}: 
    \begin{enumerate}
        \item Set $AB = s$ (or $s = 1$ for simplicity)
        \item Compute extended lengths: $AB' = 4s$, $BC' = 4s$, $CA' = 4s$
        \item Find $B'C'$ using Law of Cosines with $\angle B'BC' = 120^{\circ}$
        \item Compute area ratio: $\frac{[A'B'C']}{[ABC]} = \left(\frac{B'C'}{BC}\right)^2$
    \end{enumerate}
\end{itemize}

\subsection*{D. Conflict Resolution}
\begin{itemize}[leftmargin=*]
    \item \textbf{Calculation Check}: If Law of Cosines gives non-integer under radical, verify arithmetic
    \item \textbf{Alternative Verification}: Use area decomposition method to confirm
    \item \textbf{Answer Compatibility}: Result must match one of five answer choices exactly
\end{itemize}
\end{techniquesbox}

\begin{verificationbox}
\subsection*{A. Verification Level Selection}
\begin{itemize}[leftmargin=*]
    \item \textbf{Selected Level}: Level 3 (Standard) - 1 minute
    \item \textbf{Decision Factors}:
    \begin{itemize}
        \item Mid-band problem (Problem 15)
        \item Multiple calculation steps
        \item Answer is one of five specific choices (provides confirmation)
    \end{itemize}
\end{itemize}

\subsection*{B. Specialized Verification Methods}
\begin{itemize}[leftmargin=*]
    \item \textbf{Algebraic Check}: Verify $(3s)^2 + (4s)^2 - 2(3s)(4s)\cos(120^{\circ}) = 37s^2$
    \item \textbf{Geometric Sanity Check}: 
    \begin{itemize}
        \item Ratio must be $> 9:1$ (since extensions go beyond)
        \item Ratio should be close to but greater than $36:1$
        \item Answer (E) $37 = 36 + 1$ aligns with decomposition intuition
    \end{itemize}
    \item \textbf{Alternative Method}: Calculate using area decomposition and verify same result
    \item \textbf{Limiting Cases}: If extensions were 0, ratio would be 1:1; our ratio should grow
\end{itemize}

\subsection*{C. Confidence Calibration}
\begin{itemize}[leftmargin=*]
    \item \textbf{High Confidence Indicators}:
    \begin{itemize}
        \item Calculation yields clean answer matching choice (E)
        \item Decomposition method confirms result
        \item Answer $37 = 36 + 1$ has intuitive geometric meaning
    \end{itemize}
    \item \textbf{Target Confidence}: 95\% after verification
\end{itemize}
\end{verificationbox}

\begin{solutionbox}
\subsection*{Complete Solution}

\textbf{Method 1: Law of Cosines (Recommended)}

\textbf{Step 1: Setup and Normalization} (30 seconds)

WLOG, let $AB = BC = CA = s$. By the problem:
\begin{itemize}
    \item $BB' = 3AB = 3s$, so $AB' = AB + BB' = 4s$
    \item $CC' = 3BC = 3s$, so $BC' = BC + CC' = 4s$
    \item $AA' = 3CA = 3s$, so $CA' = CA + AA' = 4s$
\end{itemize}

By symmetry, $\triangle A'B'C'$ is also equilateral.

\textbf{Step 2: Find Angle} (15 seconds)

Since $\triangle ABC$ is equilateral, $\angle ABC = 60^{\circ}$.

Points $B$, $B'$ are collinear with $A$, so $\angle B'BC = 180^{\circ} - \angle ABC = 180^{\circ} - 60^{\circ} = 120^{\circ}$.

\textbf{Step 3: Apply Law of Cosines} (60 seconds)

For $\triangle B'BC'$, we have $BB' = 3s$, $BC' = 4s$, and $\angle B'BC' = 120^{\circ}$.

Using Law of Cosines:
\begin{align*}
    (B'C')^2 &= (BB')^2 + (BC')^2 - 2(BB')(BC')\cos(120^{\circ}) \\
    &= (3s)^2 + (4s)^2 - 2(3s)(4s)\cos(120^{\circ}) \\
    &= 9s^2 + 16s^2 - 24s^2 \cdot \left(-\frac{1}{2}\right) \\
    &= 9s^2 + 16s^2 + 12s^2 \\
    &= 37s^2
\end{align*}

Therefore, $B'C' = s\sqrt{37}$.

\textbf{Step 4: Calculate Area Ratio} (30 seconds)

Since both triangles are equilateral, they are similar. The ratio of their areas equals the square of the ratio of corresponding sides:

\[\frac{[A'B'C']}{[ABC]} = \left(\frac{B'C'}{BC}\right)^2 = \left(\frac{s\sqrt{37}}{s}\right)^2 = \left(\sqrt{37}\right)^2 = 37\]

Therefore, the answer is $\boxed{\textbf{(E) } 37:1}$.

\vspace{0.5cm}

\textbf{Method 2: Area Decomposition (Verification)}

\textbf{Step 1: Identify Components}

$[A'B'C'] = [ABC] + [AA'B'] + [BB'C'] + [CC'A']$

\textbf{Step 2: Calculate Original Triangle Area}

For equilateral triangle with side $s$: $[ABC] = \frac{s^2\sqrt{3}}{4}$

\textbf{Step 3: Calculate Corner Triangle Area}

For $\triangle AA'B'$:
\begin{itemize}
    \item $AA' = 3s$, $AB' = 4s$
    \item $\angle A'AB' = 180^{\circ} - 60^{\circ} = 120^{\circ}$
\end{itemize}

Using $A = \frac{1}{2}ab\sin\theta$:
\[[AA'B'] = \frac{1}{2} \cdot 3s \cdot 4s \cdot \sin(120^{\circ}) = \frac{1}{2} \cdot 12s^2 \cdot \frac{\sqrt{3}}{2} = 3s^2\sqrt{3}\]

By symmetry: $[BB'C'] = [CC'A'] = 3s^2\sqrt{3}$

\textbf{Step 4: Sum Areas}
\begin{align*}
    [A'B'C'] &= [ABC] + 3[AA'B'] \\
    &= \frac{s^2\sqrt{3}}{4} + 3 \cdot 3s^2\sqrt{3} \\
    &= \frac{s^2\sqrt{3}}{4} + 9s^2\sqrt{3} \\
    &= \frac{s^2\sqrt{3}}{4} + \frac{36s^2\sqrt{3}}{4} \\
    &= \frac{37s^2\sqrt{3}}{4}
\end{align*}

\textbf{Step 5: Calculate Ratio}
\[\frac{[A'B'C']}{[ABC]} = \frac{\frac{37s^2\sqrt{3}}{4}}{\frac{s^2\sqrt{3}}{4}} = 37\]

Confirmed: $\boxed{\textbf{(E) } 37:1}$
\end{solutionbox}

\begin{competitionbox}
\subsection*{A. Time Management}
\begin{itemize}[leftmargin=*]
    \item \textbf{Target Time}: 3-4 minutes for Problem 15
    \item \textbf{Time Breakdown}:
    \begin{itemize}
        \item Reading and diagram: 30 seconds
        \item Setup and angle finding: 45 seconds
        \item Law of Cosines calculation: 90 seconds
        \item Area ratio computation: 30 seconds
        \item Verification: 30 seconds
    \end{itemize}
    \item \textbf{Actual Expectation}: 3.5 minutes total
\end{itemize}

\subsection*{B. Problem Prioritization}
\begin{itemize}[leftmargin=*]
    \item \textbf{Accessibility}: High - clear geometric setup, standard techniques
    \item \textbf{Value}: 6 points (standard AMC12 scoring)
    \item \textbf{Risk Assessment}: Low risk if Law of Cosines is well-practiced
    \item \textbf{Strategic Position}: Should attempt in first pass through problems
\end{itemize}

\subsection*{C. Risk Management}
\begin{itemize}[leftmargin=*]
    \item \textbf{Confidence Threshold}: Should achieve 90\%+ confidence
    \item \textbf{Abandonment Criteria}:
    \begin{itemize}
        \item If stuck after 5 minutes, flag and move on
        \item If arithmetic becomes messy, switch to area decomposition
        \item If neither method works in 7 minutes, guess (E) based on pattern
    \end{itemize}
    \item \textbf{Error Recovery}: If calculation error, re-check angle and cosine value
\end{itemize}
\end{competitionbox}

\begin{keyinsightbox}
The key insight is recognizing that by symmetry, $\triangle A'B'C'$ is also equilateral, allowing us to find the area ratio by computing just one side length. The angle $\angle B'BC' = 120^{\circ}$ is crucial for the Law of Cosines application. The answer $37 = 36 + 1$ reflects the decomposition: three corner triangles (each 12 times the original) plus the original triangle itself.
\end{keyinsightbox}

\begin{warningbox}
\textbf{Common Pitfalls:}
\begin{itemize}
    \item Forgetting that extensions create $120^{\circ}$ angles, not $60^{\circ}$
    \item Miscalculating $\cos(120^{\circ}) = -\frac{1}{2}$ (sign error is common)
    \item Using $AB + BB' = 3s$ instead of $4s$ (forgetting to add the original side)
    \item Attempting to find area directly without recognizing the equilateral property
    \item In decomposition method, forgetting to include the original triangle $[ABC]$
\end{itemize}
\end{warningbox}

\begin{tipbox}
\textbf{Competition Tips:}
\begin{itemize}
    \item Draw a clear diagram immediately - visualization is critical
    \item Set $s = 2$ (not $s = 1$) to avoid fractions in intermediate steps
    \item Recognize answer (E) stands out as $37 = 36 + 1$, hinting at the decomposition structure
    \item If pressed for time, the area decomposition using $\frac{1}{2}ab\sin\theta$ is faster for those comfortable with it
    \item Remember: $\cos(120^{\circ}) = -\frac{1}{2}$ and $\sin(120^{\circ}) = \frac{\sqrt{3}}{2}$
\end{itemize}
\end{tipbox}

\section*{Final Answer}
\[\boxed{\textbf{(E) } 37:1}\]

\section*{Confidence Assessment}
\begin{itemize}[leftmargin=*]
    \item \textbf{Confidence Level}: 95\% - calculation verified by two independent methods
    \item \textbf{Verification Applied}: Law of Cosines (primary) + Area Decomposition (confirmation)
    \item \textbf{Flag for Review}: No - high confidence, clean calculation
    \item \textbf{Time Spent}: Approximately 3.5 minutes (on target for Problem 15)
\end{itemize}

\end{document}